\documentclass{article}
\usepackage[utf8]{inputenc}
\usepackage[T1]{fontenc}
\usepackage[margin=2cm]{geometry}
\usepackage{hyperref}
\sloppy
\begin{document}
\section*{Historical source}
Encyclopaedia Britannica scan, source row 403. Parse 3.5 extracted the text. Super 3.5 VL corrected it from an enlarged crop of each block. It can still contain recognition errors.
Page headers, footers, and page numbers are removed.

\subsection*{PHILLOOGY (E0021)}
\textit{Continues an entry from an earlier page.}

Cushim was the same. We may therefore rest assured, that whatever changes the Ethiopian dialect may have undergone in the course of three thousand years, it was originally either Chaldaic, or at least a branch of that language. Scaliger informs us, that the Ethiopians call themselves Chaldeans; and this, says he, not without reason, because of the many sacred and profane books which are extant amongst them, the most elegant and most beautiful are written in a style resembling that of the Chaldean or Assyrian. Marianus Victorius, who was the first that reduced the Ethiopic tongue to the rules of grammar, tells us, in his Prooemium, that the Ethiopians call their tongue Chaldaic; that it springs from the Babylonian, and is very like the Hebrew, Syriac, and Arabic; and at the same time, he concludes, that this language may be easily learned by those who are masters of the Hebrew. The learned Bochart, and Bishop Walton in his Prolegomena, are clearly of the same opinion. The vulgar letters of the Ethiopians, according to Diodorus Siculus, were the same with the sacred characters of the Egyptians. From this account, if the Sicilian may be trusted, the sacred letters of these people, concerning which so many vague conjectures have been formed, were actually Chaldaic. To carry on this investigation a little farther, we may observe, that Sir William Jones seems to have proved, by very plausible arguments, that the Sanscrit characters were deduced from the Chaldaic. This circumstance affords a presumption that the Ethiopian Cushim were likewise concerned with the Egyptians; who, as is remarked in the section concerning the Sancrit, probably introduced the religion of the Brahmans into Hindostan. This is advanced as a conjecture only; and yet when we consider the affinity between the Egyptian and the Gentoo religions, we

The original Ethiopians were a highly civilized people; and their laws, their institutions, and especially their religion, were celebrated far and wide. Homer talks in raptures of the piety of the Ethiopians, and sends his gods every now and then to revel twelve days with that devout people. The Sicilian produces a number of very specious arguments to prove that these two nations originally sprung from the same stock. He mentions a similarity of features, of manners, of customs, of laws, of letters, of the fabrication of statues, and of religion, as evidences of the relation between these two neighbouring nations. There was, as every body knows, a communion, in regard to sacred rites, between the two countries. The Egyptians sent annually a deputation of their priests, furnished with the portable statues of their gods, to visit the fanes of the devout Ethiopians. Upon this occasion, a solemn religious banquet was prepared, which lasted twelve days, and of which the priests of both nations were partakers. It was a kind of sacramental institution, we imagine, by which both parties publicly avouched their agreement in the ceremonies of their religion respectively. These observations plainly show, that the most ancient Ethiopians were a people highly civilized; indeed so much, that the Egyptians were at one time contented to be their scholars. The tone of their language was certainly the same with that of the Chaldeans or Arabian Cusheim, from whom they are descended. We know not whether there are any books in the ancient Ethiopian now extant; so that it is not easy to produce instances of its coincidence with the Chaldaic. But Diogenes Laertius informs us, that Thrasyllus, in his catalogue of the books composed by Democritus, mentions one,

περὶ τῶν ἐν Μερῷ ἱερῶν γραμμάτων, concerning the sacred letters in the island of Meroe; and another concerning the sacred letters in Babylon. Had these books survived the ravages of time, they would, in this age of research and curiosity, have determined not only the point under our consideration, but likewise the affinity of sacred rites amongst the Chaldeans, Ethiopians, and Egyptians.

We have now shown that the Ethiopians were a colony of Cushites; that the Cushites were originally sovereigns of Ethiopia, Shinar, or Chaldea, and consequently spoke either Chaldaic, or a dialect of that tongue; that their colonists must have used the same language; and that the ancient Ethiopians were a people highly polished, and celebrated in the most early ages on account of their virtue and piety. It has likewise appeared that the common letters of that people were the sacred characters of the Egyptians. These letters, we imagine, were the Cushic; for some account of which see the section on the Arabic. When they were discarded, and the modern letters substituted in their room, cannot be determined; nor is it, we apprehend, a matter of much importance. We shall therefore drop that part of the subject, and refer our curious and inquisitive readers to the very learned Job Ludolf's excellent grammar and dictionary of the Abyssinian or Geez tongue, where they will find every thing worth knowing on that subject. We shall endeavour to gratify our readers with a very brief account of the modern Ethiopic or Abyssinian tongue; for which we are principally indebted to that learned, indefatigable, and adventurous traveller, Bruce; who, by his observations on that country, which he made in person, often at the hazard of his life, has discovered, as it were, a new world both to Europe and Asia.

The most ancient language of Ethiopia, which we shall now call by its modern name of Abyssinia, which, according to that gentleman, the Geez, which was spoken by the ancient Cushite shepherds. This, we should think, approaches nearest to the old Chaldaic. Upon a revolution in Habesh, the court resided many years in the province of Amhara, where the people spoke a different language, or at least a very different dialect of the same language. During this interval, the Geez, or language of the shepherds, was dropped, and retained only in writing, as a dead language; the sacred Scriptures being in that tongue only saved it from going into disuse. This dialect is exceedingly harsh and unharmonious. It is full of the letters, D and T, on which an accent is put that nearly resembles stammering. Considering the small extent of sea which divides this country from Arabia, we need not wonder that it has great affinity with the Arabic. It is not difficult to be acquired by those who understand any other of the oriental languages; and as the roots of many Hebrew words are only to be found here, it seems to be absolutely necessary to all those who wish to obtain a critical skill in that language.

The Ethiopic alphabet consists of twenty-six letters, each of which, by a virgula or point annexed, varies its sound in such a manner as that those twenty-six form as it were sixty-two distinct letters. At first they had only twenty-five of these original letters, the Latin P being wanting; so that they were obliged to substitute another letter in its place. Paulus, for example, they call Taulus, Alulus, or Caulus, and Petros they pronounced Ketros. At last they substituted T, and added this to the end of their alphabet; giving it the force of P, though it was really a repetition of an old character, rather than the invention of a new one. Besides these, there are twenty others of the nature of diph-

Lib. iii. p. 101, Step
We find the same of

We find the same observation confirmed by Heliodorus (Ethiopica, lib.

We find the same observation confirmed by Heliodorus (Ethiopica, lib. x. p. 476.) "The royal letters of the Ethiopians," says he, "were the sacred characters of the Egyptians." Cassiodorus likewise assures us, "that the letters inscribed upon the Egyptian obelisks were Chaldean." But recent discoveries have refuted these and other similar statements.

Lib. ix. p. 461. Causabon.

Lib. ix. p. 461. Causabon.
Where the capital of Ethiopia was situated.
A very learned German, who published a

A very learned German, who published a grammar and dictionary of the Geez in folio.


\section*{Source record}
\begin{description}
\item[Scan ID] \texttt{britannica:309f765b1ecc83e0a7f7762e44b4a927a7555443:train:403}
\item[Image file] \texttt{page-0403-44f912e1c323.jpg}
\item[Image SHA-256] \texttt{44f912e1c3237b35750b6e188e94f3a20f2af3f1dba0a4bdb630dff640ddbf3c}
\item[Dataset] NationalLibraryOfScotland/encyclopaedia\_britannica\_illustrated at revision \texttt{309f765b1ecc83e0a7f7762e44b4a927a7555443} (CC0-1.0)
\item[Extraction] model \texttt{nvidia/nemotron-3.5-parse}, raw reply \texttt{extract/raw/page-0403-44f912e1c323.json}
\item[Raw reply SHA-256] \texttt{2cef04135ac86c240c6570ba3106b5f85bcc58d983fcea579bfcd77ce7372dee}
\item[Correction] model \texttt{nvidia/Nemotron-3.5-Super-VL}, block crops, raw replies \texttt{correct/raw/crop\_nothink/page-0403-44f912e1c323-b*.json}; super\_vl: 12
\item[Running header] PHILLOOGY.
\end{description}
\end{document}
